\documentclass[11pt]{article}

\usepackage{amsmath,amssymb,amsthm,mathrsfs}
\usepackage{geometry}
\usepackage{hyperref}
\geometry{margin=1in}

\title{
Relativistic Quantum Field Theory from Null Stochastic Geometry:\\
A Ray-Free Axiomatic Reconstruction
}

\author{}
\date{}

\theoremstyle{definition}
\newtheorem{definition}{Definition}[section]
\newtheorem{axiom}[definition]{Axiom}
\newtheorem{theorem}[definition]{Theorem}
\newtheorem{proposition}[definition]{Proposition}
\newtheorem{remark}[definition]{Remark}

\begin{document}

\maketitle

\begin{abstract}
We propose a geometric-stochastic foundation of relativistic quantum
field theory in which Hilbert space is not fundamental.
Starting from Lorentzian spacetime and a null-direction bundle endowed
with a Lorentz-invariant switching dynamics,
we reconstruct the Dirac equation, spin structure,
spin–statistics connection, second quantization,
and local operator algebras.
Renormalization appears as stochastic coarse-graining of null microstructure.
\end{abstract}

\tableofcontents

% ============================================================
\section{Introduction}

Quantum field theory is traditionally formulated via Hilbert spaces,
operator algebras, or path integrals.
We propose instead that relativistic quantum theory emerges from:

\begin{itemize}
\item Lorentzian spacetime geometry,
\item A null-direction bundle,
\item A Lorentz-invariant stochastic switching dynamics,
\item A linearization principle.
\end{itemize}

Hilbert space appears as a representation space of an underlying
geometric Markov generator.

% ============================================================
\section{Spacetime and Null Direction Bundle}

\subsection{Lorentzian Structure}

\begin{axiom}[Spacetime]
Spacetime is a smooth, oriented 4-dimensional Lorentzian manifold
$(M,g)$ of signature $(+,-,-,-)$.
\end{axiom}

\subsection{Null Direction Bundle}

\begin{definition}
The future null direction bundle is
\[
\mathcal{N}
=
\{ (x,n)\in TM \mid g(n,n)=0,\, n^0>0 \}.
\]
\end{definition}

Each fiber is diffeomorphic to $S^2$.

\begin{remark}
This bundle replaces internal spin space.
\end{remark}

% ============================================================
\section{Stochastic Worldline Dynamics}

\subsection{Worldline Process}

\begin{definition}
A particle is a Markov process
\[
(x(\tau), n(\tau)) \in \mathcal{N}
\]
with generator
\[
\mathcal{G}
=
n^\mu \nabla_\mu
+
m \mathcal{R}.
\]
\end{definition}

Here:

\begin{itemize}
\item $n^\mu \nabla_\mu$ = null transport,
\item $\mathcal{R}$ = Lorentz-invariant switching operator on $S^2$,
\item $m>0$ = switching frequency (mass parameter).
\end{itemize}

% ============================================================
\section{Clifford Algebra Emergence}

\subsection{Switching Generators}

\begin{axiom}[Clifford Structure]
There exist fiber operators $E_\mu$ satisfying
\[
\{E_\mu,E_\nu\}=2g_{\mu\nu}.
\]
\end{axiom}

\begin{theorem}
The switching algebra generates a representation of
$\mathrm{Cl}(1,3)$.
\end{theorem}

\begin{proof}[Sketch]
Lorentz invariance and bilinearity force quadratic closure
under anticommutation. Irreducibility yields the Dirac representation.
\end{proof}
% ============================================================
\section{Emergence of Gauge Interactions from Fiber Extension}

\subsection{Extended Internal Bundle}

To incorporate interactions, we enlarge the internal fiber.

\begin{definition}
Let $G$ be a compact Lie group.
Define the extended bundle
\[
\mathcal{E}
=
\mathcal{N} \times_G V,
\]
where $V$ is a unitary representation space of $G$.
\end{definition}

Each particle now carries:

\begin{itemize}
\item spacetime position $x$,
\item null direction $n$,
\item internal charge variable $q \in V$.
\end{itemize}

Thus the state space becomes
\[
(x,n,q).
\]

\subsection{Local Fiber Covariance}

Switching must remain Lorentz invariant
and now also covariant under local $G$-rotations:

\[
q(x) \mapsto U(x) q(x),
\quad U(x) \in G.
\]

\begin{axiom}[Local Covariance]
The stochastic generator must be equivariant under local $G$ transformations.
\end{axiom}

This forces introduction of a connection.

\subsection{Gauge Connection from Generator Consistency}

Replace directional derivative by covariant derivative:

\[
\nabla_\mu
\longrightarrow
D_\mu
=
\partial_\mu + A_\mu(x),
\]

where

\[
A_\mu(x) \in \mathfrak{g}
\]

is a Lie-algebra valued connection.

The worldline generator becomes

\[
\mathcal{G}
=
n^\mu D_\mu
+
m\mathcal{R}.
\]

\begin{theorem}
Local equivariance of the stochastic generator uniquely introduces
a principal $G$-connection $A_\mu$.
\end{theorem}

\begin{proof}[Sketch]
Under local transformation $U(x)$,
equivariance of $n^\mu D_\mu$ requires
\[
A_\mu \mapsto U A_\mu U^{-1}
- (\partial_\mu U)U^{-1}.
\]
Thus $A_\mu$ transforms as a gauge connection.
\end{proof}

\subsection{Dirac Equation with Gauge Coupling}

Applying the envelope principle yields

\[
(i\gamma^\mu D_\mu - m)\psi = 0.
\]

Thus minimal coupling emerges from
bundle covariance of null transport.

No gauge symmetry is postulated —
it is required for consistency of fiber switching.

% ============================================================
\subsection{Emergence of Yang–Mills Dynamics}

So far $A_\mu$ is background.
To promote it to dynamical field,
we require consistency of many-worldline interactions.

Birth–death processes involving charged particles
induce fluctuations in $A_\mu$.

Coarse-graining of worldline loops produces
an effective action for the connection:

\[
S_\text{eff}[A]
=
\int d^4x \,
\mathrm{Tr}(F_{\mu\nu}F^{\mu\nu})
+ \cdots
\]

where

\[
F_{\mu\nu}
=
\partial_\mu A_\nu
-
\partial_\nu A_\mu
+
[A_\mu,A_\nu].
\]

\begin{theorem}
One-loop worldline fluctuations generate
the Yang–Mills kinetic term.
\end{theorem}

\begin{proof}[Sketch]
Small closed zig-zag loops produce curvature contributions
via non-commutativity of parallel transport.
Expanding the worldline effective action yields $F^2$.
This reproduces the standard background-field computation.
\end{proof}

% ============================================================
\subsection{Geometric Interpretation}

Gauge field = curvature of internal fiber transport.

Mass = switching frequency.

Charge = representation of $G$ carried by internal variable $q$.

Interaction = modification of null transport by connection curvature.

Thus gauge theory is not added —
it is the geometry of extended switching bundle.

% ============================================================
\subsection{Renormalization of Gauge Couplings}

Short zig-zag loops produce scale dependence of coupling:

\[
g(\Lambda)
\]

Beta function arises from density of microscopic loops.

Thus asymptotic freedom corresponds to
decrease of effective switching coherence at short scales.

% ============================================================
\subsection{Unification Perspective}

The hierarchy now becomes:

\begin{enumerate}
\item Lorentz manifold
\item Null direction bundle
\item Extended internal bundle
\item Stochastic switching
\item Clifford algebra
\item Dirac equation with $D_\mu$
\item Worldline loop corrections
\item Yang–Mills dynamics
\item Renormalization group
\end{enumerate}

Gauge symmetry appears as
compatibility condition for local fiber transport.
% ============================================================
\section{Linear Envelope and Dirac Equation}

\subsection{Forward Kolmogorov Equation}

Probability density $\rho(x,n)$ satisfies:

\[
\partial_\tau \rho = \mathcal{G}\rho.
\]

\subsection{Linearization Principle}

\begin{axiom}[Envelope Principle]
There exists a complex envelope field $\psi(x)$
obtained from coarse-graining over fiber variables.
\end{axiom}

\begin{theorem}
The envelope satisfies the Dirac equation:
\[
(i\gamma^\mu \partial_\mu - m)\psi=0.
\]
\end{theorem}

\begin{proof}[Sketch]
Fourier analysis of null switching yields a first-order relativistic
dispersion relation.
Irreducibility of the Clifford representation fixes matrix dimension.
\end{proof}

% ============================================================
\section{Spin and Spin--Statistics}

\subsection{Spin Structure}

\begin{proposition}
Parallel transport around $2\pi$ spatial rotation induces
$\psi \mapsto -\psi$.
\end{proposition}

\subsection{Exchange Topology}

Configuration space of $N$ particles:

\[
\Gamma_N = (M\times S^2)^N / S_N.
\]

\begin{theorem}[Spin–Statistics]
Compatibility of bundle holonomy with exchange loops forces
antisymmetry of spin-$\frac12$ states.
\end{theorem}

% ============================================================
\section{Many-Worldline Extension and Fock Space}

\subsection{Configuration Space}

\[
\Gamma
=
\bigsqcup_{N=0}^\infty \Gamma_N.
\]

\subsection{Birth–Death Generator}

\[
\mathcal{G}_\text{tot}
=
\sum_i \mathcal{G}_i
+
\mathcal{B}
+
\mathcal{D}.
\]

\begin{theorem}
Linearization of the many-worldline generator produces Fock space.
\end{theorem}

\begin{proof}[Sketch]
Creation/annihilation correspond to linear representation of
configuration-space inclusion/exclusion maps.
\end{proof}

% ============================================================
\section{Renormalization as Stochastic Coarse-Graining}

\subsection{Worldline Kernel Representation}

\[
G(x,y)
=
\int_0^\infty dT \, e^{-mT} K(x,y;T).
\]

Small-$T$ divergences correspond to short zig-zag loops.

\subsection{Effective Generator}

\[
\mathcal{G}_\Lambda
=
n^\mu\partial_\mu
+
m(\Lambda)\mathcal{R}
+
\sum_i g_i(\Lambda)\mathcal{O}_i.
\]

\begin{theorem}
Running couplings arise from scale dependence of null-switching
correlation length.
\end{theorem}

% ============================================================
\section{Algebraic Quantum Field Theory Reconstruction}

\subsection{Local Algebras}

For region $\mathcal{O}\subset M$ define:

\[
\mathcal{A}(\mathcal{O})
=
\text{algebra generated by worldlines intersecting } \mathcal{O}.
\]

\subsection{Haag–Kastler Properties}

\begin{theorem}
The net $\mathcal{O}\mapsto \mathcal{A}(\mathcal{O})$
satisfies isotony and locality.
\end{theorem}

\begin{proof}[Sketch]
Null propagation enforces microcausality.
Inclusion of regions induces algebra inclusion.
\end{proof}

Vacuum = stationary measure of birth–death generator.

GNS reconstruction yields standard Fock representation.

% ============================================================
\section{Lie–Poisson Formulation}

\subsection{Bilinear Observables}

\[
Q_A = \int \bar\psi A \psi.
\]

\[
\{Q_A,Q_B\} = Q_{[A,B]}.
\]

\begin{theorem}
Dirac dynamics is Hamiltonian flow on a coadjoint orbit
of an infinite-dimensional unitary group.
\end{theorem}

% ============================================================
\section{Conceptual Hierarchy}

\begin{enumerate}
\item Lorentz manifold
\item Null direction bundle
\item Stochastic switching
\item Clifford algebra
\item Dirac equation
\item Many-worldline dynamics
\item Fock space
\item Local operator algebras
\item Renormalization group
\end{enumerate}

% ============================================================
\section{Outlook}

Future directions:

\begin{itemize}
\item Emergence of gauge symmetry from fiber extension
\item Coupling to gravity via null bundle curvature
\item Comparison with Nelson stochastic mechanics
\item Relation to worldline path integral formalism
\item Nonperturbative renormalization analysis
\end{itemize}

% ============================================================
\begin{thebibliography}{99}

\bibitem{Dirac}
P. A. M. Dirac,
\textit{The Quantum Theory of the Electron}.

\bibitem{Nelson}
E. Nelson,
\textit{Derivation of the Schrödinger Equation from Newtonian Mechanics}.

\bibitem{Haag}
R. Haag,
\textit{Local Quantum Physics}.

\bibitem{Simon}
B. Simon,
\textit{Functional Integration and Quantum Physics}.

\bibitem{DeWitt}
B. S. DeWitt,
\textit{Dynamical Theory of Groups and Fields}.

\end{thebibliography}

\end{document}
